\documentclass[11pt,a4paper]{article}

\usepackage[margin=24mm]{geometry}
\usepackage{amsmath,amssymb,bm,booktabs,tabularx,array}
\usepackage{hyperref}
\hypersetup{hidelinks}

\newcommand{\Tens}{\bm P}
\newcommand{\vect}[1]{\bm{#1}}
\newcommand{\Null}{\mathcal N}
\newcolumntype{Y}{>{\raggedright\arraybackslash}X}

\title{Tensor Polarimetry with One and Two L/R/U/D Stations}
\author{DPOLAR analysis note}
\date{}

\begin{document}
\maketitle

\section{Why a tensor-only description is appropriate}

The magnetic transport of a deuteron is governed, to leading order, by a
Hamiltonian linear in the spin,
\begin{equation}
 H(t)=-\vect\mu\cdot\vect B(t)\propto \vect B(t)\cdot\vect S.
 \label{eq:hamiltonian}
\end{equation}
The corresponding evolution operator is a spin rotation.  Under such a rotation,
irreducible spin tensors of different rank do not mix:
\begin{equation}
 U(R)\,\tau^{(k)}_q\,U^\dagger(R)
 =\sum_{q'}D^{(k)}_{q'q}(R)\tau^{(k)}_{q'}.
 \label{eq:irreducible-rotation}
\end{equation}
Vector polarization is rank one and tensor polarization is rank two.  Equation
\eqref{eq:irreducible-rotation} therefore implies that a purely tensor-polarized
beam remains purely tensor polarized during coherent magnetic transport.  The
tensor rotates, but it does not generate a vector polarization.

This separation remains true if different particles experience different
magnetic rotations, provided the trajectory weights are spin independent:
averaging rotation matrices within the beam ensemble still acts separately in
each irreducible rank.  Rank mixing would require physics
beyond Eq.~\eqref{eq:hamiltonian}, such as a relevant interaction quadratic in
spin or spin-selective loss.  Neither is part of the present beam-transport model.
It is therefore natural to restrict the polarimeter analysis to the five tensor
degrees of freedom.

In Cartesian form,
\begin{equation}
 \Tens=
 \begin{pmatrix}
 p_{xx}&p_{xy}&p_{xz}\\
 p_{xy}&p_{yy}&p_{yz}\\
 p_{xz}&p_{yz}&p_{zz}
 \end{pmatrix},\qquad
 p_{xx}+p_{yy}+p_{zz}=0.
 \label{eq:tensor}
\end{equation}
A convenient physical coordinate set is
\begin{equation}
 \vect r=(p_{xx}-p_{yy},\ p_{xy},\ p_{xz},\ p_{yz},\ p_{zz})^{\mathsf T}.
 \label{eq:physical-basis}
\end{equation}

\section{Information from one four-arm station}

With vector polarization omitted, the polarized $d$--$p$ cross section is
\begin{align}
 \frac{\sigma(\theta,\phi)}{\sigma_0(\theta)}=1
 &+\frac23(p_{xz}\cos\phi-p_{yz}\sin\phi)A_{xz}(\theta)\notag\\
 &+\frac16\left[(p_{xx}-p_{yy})\cos2\phi
 -2p_{xy}\sin2\phi\right][A_{xx}(\theta)-A_{yy}(\theta)]\notag\\
 &+\frac12p_{zz}A_{zz}(\theta).
 \label{eq:cross-section}
\end{align}
For ideal arms at
$\phi=(0^\circ,90^\circ,180^\circ,270^\circ)$, define
$B=A_{xx}-A_{yy}$.  Apart from the unpolarized rate and a common acceptance
factor, the tensor response is
\begin{equation}
 H_{\rm ring}=
 \begin{pmatrix}
 B/6&0& 2A_{xz}/3&0&A_{zz}/2\\
 -B/6&0&0&-2A_{xz}/3&A_{zz}/2\\
 B/6&0&-2A_{xz}/3&0&A_{zz}/2\\
 -B/6&0&0&2A_{xz}/3&A_{zz}/2
 \end{pmatrix}.
 \label{eq:ring-response}
\end{equation}
The measured combinations are transparent:
\begin{align}
 L-R&\propto p_{xz}A_{xz},\\
 U-D&\propto-p_{yz}A_{xz},\\
 (L+R)-(U+D)&\propto(p_{xx}-p_{yy})B.
\end{align}

The component $p_{xy}$ is absent because $\sin2\phi=0$ at all four cardinal
azimuths.  It is an exact geometrical null direction, not a weak sensitivity.
Symmetric finite detector openings do not change this conclusion.

The $p_{zz}$ term is common to all four arms.  At a single polar angle it is
indistinguishable from luminosity if the absolute normalization is free.  The
rank is then three: one diagonal anisotropy and two off-diagonal components are
measured.  If the absolute normalization is known, $p_{zz}$ supplies a fourth
independent observable.

A practical station contains more than one polar-angle group.  With one shared
luminosity, $p_{zz}$ is separated from normalization when the
acceptance-averaged $A_{zz}$ values differ between angular groups.  The station
then determines
\begin{equation}
 p_{xx}-p_{yy},\qquad p_{xz},\qquad p_{yz},\qquad p_{zz},
\end{equation}
and has rank four out of five.  The remaining null is exactly $p_{xy}$.  If each
angle has an independent unconstrained normalization, the $p_{zz}$ information
is lost and the effective rank returns to three.

For the planned two-angle proton geometry, the numerical response calculation
gives rank $4/5$ after profiling one shared luminosity.  Its null singular vector
is pure $p_{xy}$ to numerical precision.

\section{Information from two stations}

Let $\vect q_0$ denote the tensor at a common reference location.  Spin transport
to station $s$ gives
\begin{equation}
 \Tens_s=R_s\Tens_0R_s^{\mathsf T},\qquad
 \vect q_s=T^{(2)}(R_s)\vect q_0.
\end{equation}
The combined response of two stations is
\begin{equation}
 J_{\rm combined}=
 \begin{pmatrix}
 J_1T^{(2)}(R_1)\\
 J_2T^{(2)}(R_2)
 \end{pmatrix},
\end{equation}
and the unobservable tensor subspace is
\begin{equation}
 \boxed{
 \Null_{\rm combined}
 =\ker[J_1T^{(2)}(R_1)]\cap
  \ker[J_2T^{(2)}(R_2)] }.
 \label{eq:null-intersection}
\end{equation}

This equation contains the whole argument for a second polarimeter.
\begin{itemize}
 \item If the stations have the same response and the transport rotations are
 aligned, both have the same $p_{xy}$ null.  The combined rank remains four.
 The second station improves statistical precision but determines no new tensor
 component.
 \item If the relative spin rotation maps the first station's $p_{xy}$ null into
 $p_{xx}-p_{yy}$, $p_{xz}$, $p_{yz}$, or $p_{zz}$ at the second station, the
 intersection in Eq.~\eqref{eq:null-intersection} is empty.  The combined tensor
 response has full rank five.
 \item Special rotations can leave the null spaces aligned.  For example, under
 a rotation by $\alpha$ about the beam axis,
 \begin{equation}
 p'_{xy}=\frac12\sin2\alpha\,(p_{xx}-p_{yy})
          +\cos2\alpha\,p_{xy}.
 \end{equation}
 Generic angles mix $p_{xy}$ into the visible diagonal anisotropy.  At
 $\alpha=0^\circ,90^\circ,180^\circ$, the $p_{xy}$ direction maps back onto
 itself and no rank is gained.
\end{itemize}

The repository calculation shows the same behavior.  Two identical stations
with identity transport remain rank four.  Generic relative rotations about
$x$, $y$, or $z$ give rank five, except at symmetry angles where the transported
nulls realign.  These rotation scans establish the mathematical condition; an
actual beam-line conclusion requires the real spin-transfer map.

\section{Different trajectories through the beam line}

The rotation $R_s$ is an idealization.  A particle with phase-space coordinate
$\xi$ follows its own orbit and sees a rotation $R_s(\xi)$.  The tensor delivered
to the station is the ensemble average
\begin{equation}
 \overline{\Tens}_s
 =\int d\xi\,f(\xi)R_s(\xi)\Tens_0R_s^{\mathsf T}(\xi)
 =\overline{T}^{(2)}_s\vect q_0.
 \label{eq:ensemble-map}
\end{equation}
The averaged map $\overline{T}^{(2)}_s$ acts only within tensor space, so the
tensor-only description remains closed.  It is not generally an orthogonal
rotation, however.  A spread of spin phases can reduce the tensor norm and
attenuate different tensor harmonics by different amounts.

For a Gaussian spread $\sigma_\psi$ of rotations about the beam axis,
\begin{equation}
 \langle e^{im\psi}\rangle
 =e^{im\langle\psi\rangle}e^{-m^2\sigma_\psi^2/2}.
\end{equation}
The $m=1$ tensor components $(p_{xz},p_{yz})$ are reduced by
$e^{-\sigma_\psi^2/2}$, while the $m=2$ pair
$(p_{xx}-p_{yy},p_{xy})$ is reduced by $e^{-2\sigma_\psi^2}$.  The component
$p_{zz}$ is unchanged by a rotation about $z$.  An observed loss of tensor
polarization can therefore result from phase averaging even when every particle
undergoes coherent spin motion.

The detector acceptance adds a second effect.  Beam position, angle, and momentum
change the local scattering angles, and the accepted trajectories need not have
the same spin-phase distribution in L, R, U, and D.  The correct channel mean is
then
\begin{equation}
 \mu_k=\int d\xi\,f(\xi)w_k(\xi)\,
 \sigma_k\!\left[R_s(\xi)\Tens_0R_s^{\mathsf T}(\xi)\right],
 \label{eq:channel-ensemble}
\end{equation}
where $w_k(\xi)$ contains orbit acceptance and detector efficiency.  A
channel-dependent phase-space weight can imitate a tensor asymmetry or weaken the
nominal rank-five complementarity of two stations.

A quantitative two-station analysis therefore needs the orbit and spin transfer
map over the measured beam phase space, beam-position and angle information at
both stations, and finite-acceptance detector response.  Without these inputs,
two polarimeters constrain only an effective averaged tensor; they do not
separate the initial tensor from an arbitrary distribution of spin phases.

\section{Conclusion}

For magnetic transport generated by $\vect B\cdot\vect S$, vector and tensor
polarizations evolve independently.  A tensor-only analysis is therefore a
closed and physically appropriate description of the present beam.

One multi-angle cardinal L/R/U/D station determines four of the five tensor
components when a common normalization is controlled.  Its exact null direction
is $p_{xy}$.  A second aligned station improves precision but not rank.  A second
station after a known complementary spin rotation can remove the $p_{xy}$ null
and determine the complete tensor polarization.  Path-dependent spin rotations
do not create vector polarization, but their ensemble average can attenuate the
tensor and introduce additional transport and acceptance uncertainties.  The
case for two stations is consequently a statement about the transported
null-space intersection and the beam's phase-space-dependent spin map.

\end{document}
